% Build from the repository root:
% latexmk -pdf -cd overleaf/gd_imf_real_vs_calculated_slides.tex
% Redraw the figures with overleaf/scripts/build_real_vs_calculated_slides.py.
\documentclass[aspectratio=169,14pt]{beamer}
\geometry{paperwidth=254mm,paperheight=142.875mm}
% Use Beamer's standard Computer Modern fonts for a portable pdfLaTeX build.
\usepackage{fix-cm}
\DeclareFontShape{OT1}{cmss}{m}{it}{<->ssub*cmss/m/sl}{}
\usepackage{amsmath,amssymb,booktabs,array,graphicx}

\definecolor{ink}{HTML}{172B3A}
\definecolor{blue}{HTML}{2563A6}
\definecolor{orange}{HTML}{C65B32}
\definecolor{teal}{HTML}{087F8C}
\definecolor{muted}{HTML}{576B79}
\setbeamercolor{normal text}{fg=ink,bg=white}
\setbeamercolor{frametitle}{fg=ink}
\setbeamercolor{title}{fg=ink}
\setbeamercolor{structure}{fg=blue}
\setbeamercolor{alerted text}{fg=orange}
\setbeamerfont{frametitle}{size=\fontsize{25}{29}\selectfont,series=\bfseries}
\setbeamerfont{title}{size=\fontsize{35}{40}\selectfont,series=\bfseries}
\setbeamersize{text margin left=15mm,text margin right=15mm}
\setbeamertemplate{navigation symbols}{}
\setbeamertemplate{itemize items}[circle]
\setbeamertemplate{frametitle}{\vspace{6mm}\insertframetitle\par\vspace{3mm}}
\setbeamertemplate{footline}{%
    \hfill{\color{muted}\fontsize{9}{11}\selectfont\insertframenumber}\hspace{15mm}\vspace{4mm}
}
\hypersetup{colorlinks=true,urlcolor=blue,linkcolor=blue,
    pdftitle={Robust IMF under contamination},
    pdfsubject={Paired linear and robust IMF comparisons using recursive RMSE}}
\pdfstringdefDisableCommands{\def\translate#1{#1}}
\graphicspath{{figures/real_vs_calculated_slides/}}
\newcommand{\RMSE}{\operatorname{RMSE}}
\newcommand{\source}[1]{\par\vspace{2mm}{\fontsize{9}{11}\selectfont\color{muted}#1\par}}
\newcommand{\takeaway}[1]{\par\vspace{2mm}{\color{teal}#1\par}}
\newcommand{\plotwide}[2][74mm]{\par\centering\includegraphics[width=\textwidth,height=#1,keepaspectratio]{#2}\par}
\newcommand{\review}{\href{https://github.com/mkuziuk/imf/blob/main/overleaf/gd_imf_real_vs_calculated_review.tex}{Source review}}
\newcommand{\notebooklink}[2]{\href{https://github.com/mkuziuk/imf/blob/main/experiments/real-vs-calculated/#1}{#2}}

\title{Robust IMF under contamination}
\author{}
\date{}

\begin{document}

\begin{frame}[plain]
    \vfill
    {\usebeamerfont{title}Robust IMF under contamination\par}
    \vfill
\end{frame}

\begin{frame}{Robust IMF changes the local fitting problem}
    Set \(Y^{(1)}=Y\). At each stage \(k\), extract the IMF component
    \[
        \widetilde S^{(k)}_t=\operatorname*{arginf}_{x}\;
        \sum_u \rho\!\left(Y^{(k)}_u-x\right)K_{h_k}(t-u).
    \]
    \(K_h(v)=K(v/h)\) weights the local window. The contrast \(\rho\) penalizes
    \(r=Y^{(k)}_u-x\); its derivative \(\psi=\rho'\) determines each point's contribution to the fit.
    \vspace{4mm}
    \begin{columns}[T,onlytextwidth]
        \begin{column}{0.46\textwidth}
            {\color{blue}\large Linear IMF}\par
            \[
                \rho_{\rm lin}(r)=\frac{r^2}{2},\qquad \psi_{\rm lin}(r)=r.
            \]
            A weighted mean; large mismatches have large influence.
        \end{column}
        \begin{column}{0.49\textwidth}
            {\color{orange}\large Robust IMF}\par
            \[
                \rho_H(r)=\mathbb E_{\xi}\,|r+H\xi|,\qquad \xi\sim\mathcal N(0,1),
            \]
            \[
                \psi_H(r)=\operatorname{erf}\!\left(\frac{r}{\sqrt{2}H}\right),\qquad |\psi_H(r)|\leq1.
            \]
            A smooth absolute-value contrast, fitted by gradient descent; \(H=2\sigma\).
        \end{column}
    \end{columns}
    \vspace{4mm}
    Subtract the component and shrink the bandwidth:
    \[
        Y^{(k+1)}_t=Y^{(k)}_t-\widetilde S^{(k)}_t,
        \qquad h_{k+1}=h_k/a,\qquad a=\sqrt{2}.
    \]
\end{frame}

\begin{frame}{Compare both methods on the same observations}
    \begin{columns}[T,onlytextwidth]
        \begin{column}{0.46\textwidth}
            {\color{blue}\large Additive contamination}
            \[
                Y_i=X_i+G_i+M_iC_i.
            \]
            Contaminated points retain the signal.
        \end{column}
        \begin{column}{0.49\textwidth}
            {\color{blue}\large Masked pure-noise contamination}
            \[
                Y_i=\begin{cases}
                    X_i+G_i,&M_i=0,\\
                    G_i+C_i,&M_i=1.
                \end{cases}
            \]
            Contaminated points lose the signal.
        \end{column}
    \end{columns}
    \vspace{4mm}
    \(G_i\sim\mathcal N(0,\sigma^2)\), \(M_i\sim\operatorname{Bernoulli}(p)\),
    and \(C_i\) is exponential noise with an independent random sign.\par
    \vspace{4mm}
    The RMSE curves and heatmaps compare each method \(m\) with its matching clean reference:
    \[
        e_{k,m}(t_i)=\widetilde S^{(k)}_m(Y)(t_i)-\widetilde S^{(k)}_m(X)(t_i),
        \qquad \RMSE_{k,m}=\sqrt{\frac{1}{n}\sum_{i=1}^{n} e_{k,m}(t_i)^2}.
    \]
    \takeaway{Both methods receive the same observation. Stage-comparison plots cover all nine stages.}
\end{frame}

\begin{frame}{The contaminated observations}
    \plotwide[71mm]{paired_observations.pdf}
    \raggedright
    \(n=1000\), \(\sigma=0.6\), \(p=0.2\), contamination scale \(2.0\), robust \(H=1.2\).\par
    Nine windows shrink from 501 to 31 points. The 175 contaminated points are marked in orange.
\end{frame}

\begin{frame}{Additive contamination at the default setting}
    \plotwide{additive_method_rmse.pdf}
    \raggedright
    Mean RMSE across the nine stages falls from \(0.0611\) to \(0.0394\), a \textcolor{teal}{35.5\% reduction}.\par
    Robust IMF has lower RMSE at all nine stages in this realization.
\end{frame}

\begin{frame}{Masked contamination at the default setting}
    \plotwide{masked_method_rmse.pdf}
    \raggedright
    Mean RMSE falls from \(0.0654\) to \(0.0437\), a \textcolor{teal}{33.2\% reduction}.\par
    The robust advantage appears at every stage in this paired realization.
\end{frame}

\begin{frame}{Robust IMF improves the first masked component}
    \plotwide{masked_first_stage_methods.pdf}
    \raggedright
    Both curves use the same clean linear IMF component as their reference.\par
    Stage-1 RMSE is \(0.1234\) for linear IMF and \(0.0898\) for robust IMF,
    a \textcolor{teal}{27.3\% reduction}.
\end{frame}

\begin{frame}{The advantage grows with contamination magnitude}
    \plotwide{contamination_magnitude_rmse.pdf}
    \raggedright
    Fixed \(\sigma=0.2\), \(p=0.2\), \(H=0.4\); curves and bands show medians and interquartile ranges across three repeats.\par
    At scale \(12\sigma\), robust mean RMSE is \textcolor{teal}{74.6\% lower} in the additive model
    and \textcolor{teal}{77.1\% lower} in the masked model.
\end{frame}

\begin{frame}{The contamination grid compares both methods}
    \plotwide[80mm]{sweep_heatmaps.pdf}
    \raggedright
    Additive contamination; three repeats per cell. Stars mark the default setting.\par
    Both heatmaps use one logarithmic color scale for the median mean RMSE over nine stages.
\end{frame}

\begin{frame}{The masked grid compares both methods}
    \plotwide[80mm]{masked_sweep_heatmaps.pdf}
    \raggedright
    Masked contamination; three repeats per cell. Stars mark the default setting.\par
    Both heatmaps use one logarithmic color scale for the median mean RMSE over nine stages.
\end{frame}

\begin{frame}{Robust gains across all nine stages at \(12\sigma\)}
    \plotwide{strong_contamination_stage_rmse.pdf}
    \raggedright
    At \(\sigma=0.2\), \(p=0.2\) and contamination scale \(12\sigma=2.4\), median RMSE over three repeats is lower\par
    at all nine stages in both models. Stage-wise reductions range from \textcolor{teal}{73\% to 80\%}.
\end{frame}

\begin{frame}{Larger masked replacements can reduce robust RMSE}
    \plotwide{masked_magnitude_rmse.pdf}
    \raggedright
    Robust mean RMSE falls \textcolor{teal}{13\%} from \(2\sigma\) to \(8\sigma\), then rises; linear RMSE increases.\par
    Robust stage-1 RMSE falls \textcolor{teal}{56\%} from \(2\sigma\) to \(12\sigma\), while linear RMSE rises \textcolor{blue}{70\%}.\par
    Fixed \(\sigma=0.2\), \(p=0.2\), \(H=0.4\); medians and interquartile ranges across three repeats.
\end{frame}

\begin{frame}{Why larger replacements may be less harmful}
    At a masked point, \(Y_i=G_i+C_i\). The clean signal is absent at every contamination magnitude.
    \vspace{6mm}
    \begin{columns}[T,onlytextwidth]
        \begin{column}{0.46\textwidth}
            {\color{blue}\large Small replacement values}\par\vspace{3mm}
            Replacements concentrate near zero. They can pull a local fit toward zero where the clean signal is nonzero.
        \end{column}
        \begin{column}{0.49\textwidth}
            {\color{orange}\large Large, symmetric replacements}\par\vspace{3mm}
            The bounded robust score limits the influence of large residuals. Positive and negative score contributions can cancel.
        \end{column}
    \end{columns}
    \vspace{5mm}
    \[
        |\psi_H(r)|\leq1,
        \qquad \frac{\psi_H(r)}{r}\longrightarrow0\quad\text{as }|r|\to\infty.
    \]
    \takeaway{The observed RMSE curve is consistent with this explanation. The experiment does not isolate the mechanism.}
\end{frame}

\begin{frame}{The numerical case for robust IMF}
    \begin{itemize}
        \setlength{\itemsep}{7mm}
        \item Lower RMSE at all nine stages in both single-run examples.\par
            Mean RMSE is 35.5\% lower for additive contamination and 33.2\% lower for masked contamination.
        \item Much larger gains when contamination values are large.\par
            At scale \(12\sigma\), mean RMSE falls by about 75\% and 77\% in the two models.
        \item The advantage under large contamination spans the full decomposition.\par
            Median RMSE improves at every stage, with reductions of 73\% to 80\%.
    \end{itemize}
    \vspace{5mm}
    \takeaway{A bounded local score preserves the IMF procedure while reducing sensitivity to large contamination values.}
\end{frame}

\appendix

\begin{frame}{Appendix: implementation details}
    \begin{itemize}
        \setlength{\itemsep}{5mm}
        \item Window sizes are \(501,355,251,177,125,89,63,45,31\).\par
            Both methods use periodic wrap boundaries and the same weights.
        \item The implemented weights are \(w(u)\propto(1-|u|)_+^2\).\par
            The review calls the kernel Epanechnikov; that label differs from the code.
        \item Robust fits use \(H=2\sigma\), the notebook's lookup grid,\par
            at most 60 gradient-descent iterations and tolerance \(10^{-6}\).
        \item RMSE curves and heatmaps use each method's clean reference.\par
            The waveform slide uses one shared clean linear IMF reference.
    \end{itemize}
    \source{Main example: \(\sigma=0.6\), \(p=0.2\), contamination scale \(2.0\); signal seed 2026, observation seed 777.
    Sweep: \(\sigma=0.2\), \(p=0.2\), scales \(2\sigma,4\sigma,\ldots,12\sigma\), seeds 777000--777002.}
\end{frame}

\begin{frame}{Appendix: sources and reproducibility}
    \begin{itemize}
        \setlength{\itemsep}{5mm}
        \item \texttt{IMF.pdf}, \S2.2 and equation (3.2)\par
            Component definition and recursive update; robust contrast in \S3.
        \item \review\par
            Source experiment, contamination models and robust contrast.
        \item \notebooklink{gd_imf_real_vs_calculated.ipynb}{Additive comparison} and
            \notebooklink{gd_imf_real_vs_calculated_noise_only.ipynb}{masked comparison}\par
            Both methods are recomputed on the same contaminated observation.
        \item \notebooklink{gd_imf_real_vs_calculated_limits.ipynb}{Additive sweep} and
            \notebooklink{gd_imf_real_vs_calculated_limits_noise_only.ipynb}{masked sweep}\par
            Magnitude slices: 18 trials per model. Heatmaps: 225 trials per model across three noise levels,
            five probabilities and five magnitudes, with three repeats per setting.
    \end{itemize}
    \vspace{5mm}
    The companion script exports stage RMSE, mean RMSE and vector figures.\par
    CSV data and source hashes are stored beside the figures.
    \source{Build instructions: \texttt{gd\_imf\_real\_vs\_calculated\_slides.md}, next to the LaTeX source.}
\end{frame}

\end{document}
